\section{Military food supply-chain case and experimental genealogy}\label{sec:case}

The case is a Python/SimPy reconstruction of the 13-operation military food supply chain
reported by Garrido-Rios~\cite{garrido2017}. It transforms raw materials into assembled
rations and distributes them through a Supply Battalion and downstream units under production,
transport, demand, and disruption risks. The reconstruction reproduces the workbook-level ReT
formula and retains order ledgers, processing times, queues, resource use, and disruption
tapes. We describe it as thesis-grounded reconstruction with forensic workbook replay and
throughput checks, not as a validated digital twin. Extensions such as finite convoys or split
CSSUs are disclosed as researcher-imposed decision contracts.

The empirical program proceeds from cheap causal diagnostics to learning. Track A evaluates
thesis-native buffers and shifts. Track B adds dispatch and initially compares PPO with a
restricted static family; a later same-contract challenge grants constants the learner's full
rights. Program D then screens candidate decision families using exact replay. DRA-1 introduces
daily CSSU allocation with unused-capacity reallocation. DRA-2 and DRA-2b introduce a finite
convoy and enumerate restricted open-loop sequences. Program E is the first clean learner study
opened only after material clairvoyant convoy headroom was observed. Program F studies a fixed
readiness budget allocated among maintenance, transport protection, and reserve. Program G
studies spatial priority between two stylized CSSUs under persistent tempo and imperfect
signals. Every program has its own seed universe and terminal verdict.

Three layers of evidence must not be merged. The full DES supports claims about the reconstructed
Op1--Op13 chain and its order metric. The convoy studies are disclosed physical extensions of
that chain. Program G is a stylized weekly spatial order adapter used to isolate information and
allocation mechanisms; it is not a port of the spatial mechanism into the full event-driven
DES. Its corrected terminal audit uses a 168-hour calendar, six operating days, the canonical
cumulative ReT ledger, and unattended orders marked lost at the horizon.

All comparisons use exact or common-random-number tapes. Static policies are selected on
calibration tapes and frozen before holdout. Observable policies exclude future demand,
future disruption duration, latent regime labels, and oracle outcomes. Resource accounting
includes the relevant decision-family quantities: shifts and dispatch multipliers in Track B,
departures and unavailable convoy hours in Program D/E, two readiness tokens in Program F,
and identical spatial action rights in Program G. PPO is prohibited until the program-specific
pre-learning gates pass.

The primary continuity outcome is the workbook-faithful order-level ReT. We report service
loss, fill, backlog, unattended orders, completion time, recovery measures, and resource use as
guardrails. For the spatial audit we additionally report quantity-weighted ReT and worst-CSSU
fill. Cobb--Douglas scores use fixed published-family exponents only as secondary construct
sensitivity because their factory-planning calibration is not a calibration for this MFSC.

